A continuación, presentamos dos herramientas fundamentales para el 
estudio de la equidistribución. La primera es una caracterización funcional,
que nos permite estudiar la equidistribución a través de las medias de
funciones integrables Riemann. La segunda es el criterio de Weyl, que proporciona 
una caracterización en términos de sumas exponenciales.  

\section{Caracterización funcional}

\begin{theorem}
    \label{carac:riemann}
    Sea $(x_n)_{n\geq1}$ una sucesión de números reales, entonces son equivalentes:
    \begin{enumerate}
        \item $(x_n)$ es equidistribuida módulo $1$.
        \item Para toda función $f\in \mathcal{R}(\mathbb{T})$, Riemann integrable, se cumple que
              \[
                  \lim_{N \to \infty} \frac{1}{N} \sum_{n=1}^N f(x_n) = \int_0^1 f(x)\ dx.
              \]
    \end{enumerate}
\end{theorem}

\begin{proof}
    $1 \Longrightarrow 2$ \, Sea $(x_n)$ una sucesión equidistribuida módulo 1.
    
    Consideremos la partición $0=a_0 < a_1 < \ldots < a_k=1$ de $\mathbb{T}$. Sea $f$ una función 
    escalonada tal que $f(x) = \sum_{i=1}^k c_i \mathbf{1}_{[a_{i-1}, a_i)}(x)$ con $c_i \in \mathbb{R}$
    para cada $i=1,\ldots,k$.

    Entonces,
    \[
        \frac{1}{N} \sum_{n=1}^N f(x_n) = 
        \frac{1}{N} \sum_{n=1}^N \sum_{i=1}^k c_i \mathbf{1}_{[a_{i-1}, a_i)}(x_n) = 
        \sum_{i=1}^k c_i \left(\frac{1}{N} \sum_{n=1}^N \mathbf{1}_{[a_{i-1}, a_i)}(x_n)\right)
    \]
    como $(x_n)$ es equidistribuida, por la observación~\ref{obs:equidistribucion}, 
    \begin{align*}
        \lim_{N \to \infty} \sum_{i=1}^k c_i \left(\frac{1}{N} \sum_{n=1}^N \mathbf{1}_{[a_{i-1}, a_i)}(x_n)\right)
        &= \sum_{i=1}^k c_i \left(\lim_{N \to \infty} \frac{1}{N} \sum_{n=1}^N \mathbf{1}_{[a_{i-1}, a_i)}(x_n)\right) \\
        &= \sum_{i=1}^k c_i \int_0^1 \mathbf{1}_{[a_{i-1}, a_i)}(x)\ dx \\
        &= \int_0^1 \sum_{i=1}^k c_i \mathbf{1}_{[a_{i-1}, a_i)}(x)\ dx 
        =\int_0^1 f(x)\ dx.
    \end{align*}
    
    Ahora, sea $f \in \mathcal{R}(\mathbb{T})$. Como $f$ es integrable Riemann, 
    para todo $\varepsilon > 0$ existe una partición $P=\{0=a_0<\cdots<a_k=1\}$ tal que, 
    si definimos
    \[
        L(x)=\sum_{i=1}^k m_i \mathbf{1}_{[a_{i-1},a_i)}(x) \qquad
        U(x)=\sum_{i=1}^k M_i \mathbf{1}_{[a_{i-1},a_i)}(x)
    \]
    donde $m_i=\inf_{x\in[a_{i-1},a_i)} f(x)$ y $M_i=\sup_{x\in[a_{i-1},a_i)} f(x)$
    se tiene
    \[
        L(x)\le f(x)\le U(x), \qquad \int_0^1 U(x) \, dx - \int_0^1 L(x) \, dx \leq \varepsilon.
    \]

    Sea $S_N = \sum_{n=1}^N f(x_n)$, tenemos,
    \[
        \frac{1}{N} \sum_{n=1}^N L(x_n) \leq \frac{S_N}{N} \leq \frac{1}{N} \sum_{n=1}^N U(x_n)
    \]
    como $L$ y $U$ son funciones escalonadas, al tomar límites,
    \[
        \int_0^1 f(x) \, dx - \varepsilon 
        \leq \int_0^1 L(x) \, dx 
        \leq \liminf_{N \to \infty} \frac{S_N}{N} 
        \leq \limsup_{N \to \infty} \frac{S_N}{N} 
        \leq \int_0^1 U(x) \, dx 
        \leq \int_0^1 f(x) \, dx + \varepsilon
    \]
    como $\varepsilon > 0$ es arbitrario, se concluye que
    \[
        \lim_{N \to \infty} \frac{1}{N} \sum_{n=1}^N f(x_n) = \int_0^1 f(x) \, dx.
    \]

    $2 \Longrightarrow 1$ \, Sea un intervalo $(a,b) \subset \mathbb{T}$ y sea $f(x) = \mathbf{1}_{(a,b)}(x)$, 
    como $f$ es Riemann integrable, entonces,
    \[
        \lim_{N \to \infty} \frac{1}{N} \sum_{n=1}^N \mathbf{1}_{(a,b)}(x_n) = \int_0^1 \mathbf{1}_{(a,b)}(x)\ dx = b-a
    \]

    lo que demuestra que $(x_n)$ es equidistribuida módulo $1$.
\end{proof}

\begin{corollary}
    \label{carac:continua}
    Sea $(x_n)_{n\geq1}$ una sucesión de números reales, entonces son equivalentes:
    \begin{enumerate}
        \item $(x_n)$ es equidistribuida módulo $1$.
        \item Para toda función $f\in C(\mathbb{T})$, se cumple que
              \[
                  \lim_{N \to \infty} \frac{1}{N} \sum_{n=1}^N f(x_n) = \int_0^1 f(x)\ dx.
              \]
    \end{enumerate}
\end{corollary}

\begin{example}
    Sea $(x_n)_{n\geq1}$ una sucesión en $\mathbb{T}$. Definimos $S((a,b),N) = \sum_{n=1}^N x_n \, \mathbf{1}_{(a,b)}(x_n)$
    donde $(a,b) \subseteq \mathbb{T}$. Entonces, la sucesión $(x_n)$ es equidistribuida módulo $1$ si y solo si 
    para todo intervalo $(a,b) \subseteq \mathbb{T}$
    \[
        \lim_{N \to \infty} \frac{S((a,b),N)}{N} = \frac{1}{2}(b^2-a^2)
    \]
\end{example}

\begin{proof}
    $\Longrightarrow$ \, Sea $(a,b) \subseteq \mathbb{T}$ y sea $f(x) = x \, \mathbf{1}_{(a,b)}(x)$.
    
    Como $f \in \mathcal{R}(\mathbb{T})$ y $(x_n)$ es equidistribuida,
    \[
        \lim_{N \to \infty} \frac{S((a,b),N)}{N} = 
        \lim_{N \to \infty} \frac{1}{N} \sum_{n=1}^N f(x_n) = 
        \int_0^1 f(x)\ dx = \int_a^b x\ dx = \frac{1}{2}(b^2-a^2)
    \]

    $\Longleftarrow$ \, 
    Vamos a probar que las CL de funciones $f(x)=x \, \mathbf{1}_{(a,b)}(x)$ son densas en $L^1(\mathbb{T})$.
    Para esto, sea $f \in L^1(\mathbb{T})$ y sea $\varepsilon > 0$. Por la densidad de las funciones simples, existe una función simple $s(x) = \sum_{i=1}^m a_i \mathbf{1}_{A_i}(x)$ tal que $\|f-s\|_1 < \varepsilon/2$.
    Dado que cada $A_i$ es medible, podemos aproximar cada $A_i$ por una unión finita de intervalos abiertos, es decir, existe una función $g(x) = \sum_{i=1}^m a_i \mathbf{1}_{B_i}(x)$ donde cada $B_i$ es una unión finita de intervalos abiertos tal que $\|s-g\|_1 < \varepsilon/2$.
    Por lo tanto, $\|f-g\|_1 \leq \|f-s\|_1 + \|s-g\|_1 < \varepsilon$.
    Finalmente, cada $B_i$ es una unión finita de intervalos abiertos, por lo que $g$ es una combinación lineal de funciones de la forma $x \, \mathbf{1}_{(a,b)}(x)$.      
\end{proof}

\section{El criterio de Weyl}

La caracterización funcional de la equidistribución es una herramienta de gran
utilidad, pero a menudo resulta difícil de aplicar a los distintos problemas. 
El criterio de Weyl proporciona un método más concreto para verificar la 
equidistribución, utilizando sumas exponenciales. 

Este criterio se basa en el hecho de que las funciones $f(x)=e^{2\pi i k x}$ con 
$k\in\mathbb{Z}$ son continuas y por tanto, satisfacen las condiciones
del corolario \ref{carac:continua}. Además, esta familia de funciones es lo suficientemente 
rica como para caracterizar la equidistribución.

\begin{theorem}
    \label{teo:criterio_weyl}
    Sea $(x_n)$ una sucesión en $[0,1)$, entonces son equivalentes:
    \begin{enumerate}
        \item $(x_n)$ es equidistribuida en $[0,1)$.
        \item Para todo $k \in \mathbb{Z} \setminus \{0\}$,
              \[
                  \frac{1}{N} \sum_{n=1}^N e^{2\pi i k x_n} \longrightarrow 0 \quad \text{cuando } N \to \infty.
              \]
    \end{enumerate}
\end{theorem}

\begin{proof}
    $1 \Longrightarrow 2$ \, Sea $k \in \mathbb{Z} \setminus \{0\}$, como $(x_n)$ es equidistribuida aplicando
    el teorema anterior a la función $f(x) = e^{2\pi i k x}$, entonces,
    \[
        \lim_{N \to \infty} \frac{1}{N} \sum_{n=1}^N e^{2\pi i k x_n} = \int_0^1 e^{2\pi i k x}\ dx = 0.
    \]
    $2 \Longrightarrow 1$ \, La hipótesis es el paso $2$ del Lema \ref{lem:lema_na}. Reproduciendo la prueba,
    se concluye que $(x_n)$ es equidistribuida módulo $1$. 
\end{proof}


\begin{example}
    La sucesión $(n\alpha)_{n\geq1}$ con $\alpha \notin \mathbb{Q}$ es equidistribuida módulo $1$.
\end{example}

\begin{resolve}
    Sea $k \in \mathbb{Z}\setminus\{0\}$,
    \[
        \left|\frac{1}{N}\sum_{n=1}^N e^{2\pi i k n\alpha} \right| =
        \frac{1}{N} \left|\frac{e^{2\pi i k \alpha}(1-e^{2\pi i k N \alpha})}{1-e^{2\pi i k \alpha}} \right| \leq
        \frac{2}{N|e^{2\pi i k \alpha}-1|}
    \]
    Como $\alpha \notin \mathbb{Q}$, se tiene $k\alpha \notin \mathbb{Z}$ y por tanto $e^{2\pi i k \alpha} \neq 1$
    
    Al tomar el límite cuando $N \to \infty$ se concluye que
    \[
        \lim_{N \to \infty} \frac{1}{N}\sum_{n=1}^N e^{2\pi i k n\alpha} = 0
    \]
    y por el criterio de Weyl, la sucesión $(n\alpha)$ es equidistribuida módulo $1$.
\end{resolve}



\begin{example} 
    \label{ex:ax-b}
    Sea $(x_n)_{n\geq1}$ una sucesión equidistribuida módulo $1$, sea $a\in \mathbb{Z}\setminus\{0\}$ y sea $b\in\mathbb{R}$.
    Entonces la sucesión $(ax_n+b)_{n\geq1}$ es equidistribuida módulo $1$.
\end{example}

\begin{resolve}
    Definimos $y_n = ax_n+b$. Para cada $k \in \mathbb{Z}\setminus\{0\}$ se tiene
    \[
        \left|\frac{1}{N}\sum_{n=1}^N e^{2\pi i k y_n} \right|
        = \left|\frac{1}{N}\sum_{n=1}^N e^{2\pi i k (ax_n + b)} \right|
        = \left|\frac{1}{N}e^{2\pi i k b}\sum_{n=1}^N e^{2\pi i (ka) x_n} \right|
        =\left|\frac{1}{N}\sum_{n=1}^N e^{2\pi i (ka) x_n} \right|.
    \]
    
    Dado que $ka \in \mathbb{Z}\setminus\{0\}$ y $(x_n)$ es equidistribuida, el criterio de Weyl implica que
    \[
        \lim_{N\to\infty} \frac{1}{N}\sum_{n=1}^N e^{2\pi i (ka) x_n} = 0.
    \]
    
    Por consiguiente,
    \[
        \lim_{N\to\infty} \frac{1}{N}\sum_{n=1}^N e^{2\pi i k y_n}  = 0.
    \]
    y por el criterio de Weyl, la sucesión $(y_n)$ es equidistribuida módulo $1$.
\end{resolve}

Como consecuencia, la proposición \ref{prop:secuencia_mas_const} es un corolario de este resultado.

\begin{example}
    La sucesión $\partfrac{\sin n}$ es densa en $[0,1)$ pero no es equidistribuida.
\end{example} 

\begin{resolve}
    Sabemos que $n$ es equidistribuida en $[0,2\pi)$.
    Por la caracterización funcional de la equidistribución aplicada a la función $f(x) = e^{2\pi i \sin x}$, tenemos que
    \[
        \frac{1}{N} \sum_{n=1}^N e^{2\pi i \sin n} \longrightarrow \frac{1}{2\pi} \int_0^{2\pi} e^{2\pi i \sin x} dx \qquad \text{cuando } N \to \infty.
    \]

    Si calculamos esta integral, tenemos que
    \[
        \int_0^{2\pi} e^{2\pi i \sin x} dx = 2\pi J_0(2\pi) \neq 0.
    \]

    Si $\partfrac{\sin n}$ fuese equidistribuida, por el criterio de Weyl, debería cumplirse
    \[
        \forall \, k \in \mathbb{Z}\setminus\{0\}, \quad \frac{1}{N} \sum_{n=1}^N e^{2\pi i k \sin n} \longrightarrow 0 \qquad \text{cuando } N \to \infty,
    \]
    Pero para $k=1$ hemos visto que esta suma tiende a $J_0(2\pi) \neq 0$, lo que demuestra que $\partfrac{\sin n}$ no es equidistribuida.
\end{resolve}

